\documentclass[12pt]{article}

\usepackage[utf8]{inputenc}
\usepackage[spanish, es-noshorthands]{babel}

\usepackage{mathtools}
\usepackage{amssymb}
\usepackage{amsthm}

\usepackage{geometry}
\geometry{margin=2.5cm}

\usepackage{tikz}
\usetikzlibrary{positioning}

\usepackage{hyperref}

\title{Entregables Teoría de Números}
\author{Juan Felipe Rodríguez Córdoba}
\date{}

\begin{document}
\maketitle

\section*{Ejercicio 1}
Se sabe que
\[
3 \mid 4^n - 1.
\]
Demuestra que
\[
9 \mid 4^n - 3n - 1
\]

\section*{Ejercicio 2}
¿Es verdad que para cualquier \(n \in \mathbb{N}\) resulta que
\[
P(n)=n^2+n+41
\]
es primo?

\section*{Ejercicio 3}
¿Cuántos divisores tienen estos números?
\[
11\cdot 13,
\]
\[
11\cdot 13\cdot 17,
\]
\[
11^{14},
\]
\[
11^{17}\cdot 13^{28}\cdot 17^{65}.
\]

\section*{Ejercicio 4}
Demuestra que para cualquier primo \(p>3\) se cumple que
\[
p^2-1
\]
es divisible por \(24\).

\section*{Ejercicio 5}
Demuestra que
\[
a^3+b^3+c^3-3abc
\]
es divisible por
\[
a+b+c.
\]

\section*{Ejercicio 6}
Demuestra que para cualquier \(n\in\mathbb{N}\) se cumple que
\[
k! \mid n(n+1)(n+2)\cdots(n+k-1).
\]

\section*{Ejercicio 7}
Sea
\[
n! = p_1^{a_1}p_2^{a_2}\cdots p_s^{a_s}.
\]
Demuestra que
\[
a_k=
\left\lfloor \frac{n}{p_k}\right\rfloor+
\left\lfloor \frac{n}{p_k^2}\right\rfloor+
\left\lfloor \frac{n}{p_k^3}\right\rfloor+\cdots.
\]

\section*{Ejercicio 8}
Demuestra que para cualquier primo \(p\) se cumple que
\[
(p-1)! \equiv -1 \pmod p.
\]

\section*{Ejercicio 9}
Demuestra que \(p\) es primo si y solo si
\[
(p-2)! \equiv 1 \pmod p.
\]

\section*{Ejercicio 10}
Si \(p\) y \(q\) son primos distintos, demuestra que
\[
p^q+q^p \equiv p+q \pmod{pq}.
\]

\section*{Ejercicio 11}
Busca
\[
\gcd(2n+13,n+7).
\]

\section*{Ejercicio 12}
Busca
\[
\gcd(2^{100}-1,2^{120}-1).
\]

\section*{Ejercicio 13}
Demuestra que si
\[
\gcd(m,n)=1,
\]
entonces
\[
a\equiv b \pmod{mn}
\]
si y solo si
\[
a\equiv b \pmod m
\]
y
\[
a\equiv b \pmod n.
\]

\section*{Ejercicio 14}
Busca el mínimo número par \(a\) tal que
\[
3\mid a+1,\qquad 5\mid a+2,\qquad 7\mid a+3,
\]
\[
11\mid a+4,\qquad 13\mid a+5.
\]

\section*{Ejercicio 15}
¿Para qué números enteros \(n\) se cumple que
\[
55\mid n^2+3n+1?
\]

\section*{Ejercicio 16}
En una tabla cuadrada hemos coloreado una tabla cuadrada más pequeña.
Han quedado \(79\) casillas sin colorear.
¿Pueden quedar las esquinas sin colorear?

\section*{Ejercicio 17}
Hemos cogido todos los números del \(1\) al \(1000000\) y hemos cambiado cada número por la suma de sus cifras.
Luego hemos repetido este proceso hasta obtener una serie de números de una sola cifra.
¿Cuántos más unos que doses hay?
\end{document}

\section*{Hoja 3}

\section*{Ejercicio 1}
Representa
\[
\sqrt{7}
\]
como fracción continua periódica.

\section*{Ejercicio 2}
Demuestra que hay infinitos primos del tipo
\[
4n+3.
\]

\section*{Ejercicio 3}
Busca todos los números de \(3\) cifras cuyo cuadrado termina en este propio número.

\section*{Ejercicio 4}
Demuestra que la ecuación
\[
4x^2+3y^2=2000000002
\]
no tiene solución en números enteros.

\section*{Ejercicio 5}
Resuelve la ecuación
\[
x^2\equiv 1\pmod{35}.
\]
Justifica tu solución.

\section*{Ejercicio 6}
Demuestra que para cualquier \(N\) existe una lista de \(N\) números compuestos sucesivos.

\section*{Ejercicio 7}
Dado el mensaje cifrado
\[
\texttt{t5r](+},
\]
descifra el mensaje si para cifrarlo se utilizó el algoritmo RSA con
\[
n=77,\qquad e=37,
\]
y la equivalencia alfanumérica dada en la hoja.

\section*{Ejercicio 8}
Demuestra que para cualquier \(\varepsilon>0\) existe un número grande \(N\) tal que entre los primeros \(N\) números el porcentaje de primos es menor que \(\varepsilon\).